\chapter{Corps}

\setlength{\parindent}{0pt}
\renewcommand{\labelitemi}{\textbullet} % Utiliser des points noirs (•)

% ==================================================================================================================================
% Introduction

Résumons, nous avons un groupe avec un loi supplémentaire qui peut être commutative et posséder un neutre. 
Grâce à cela, nous avons pu définir les idéeaux, un nouveau type de morphisme, etc... 

Seulement, on souhaiterai pouvoir définir des structures plus complexes comme avec les groupes tels que les anneaux 
quotient, avoir une notion de divisibilité, bref, de l'arithmétique. 

% ==================================================================================================================================
% Corps 

\section{Corps, définition et propriétés}

\begin{definition}[Corps]
    On appelle $(K, +, \times)$ un corps si $(K,+,\times)$ est un anneau et $(K^*, \times)$ un groupe. 
    Si la loi $\times$ est commutative, on parlera alors de corps commutatif. 
\end{definition}

On peut bien évidement trouver une condition nécessaire et suffisante pour qu'un anneau soit un corps...

\begin{proposition}
    Soit $(K, +, \times)$ un anneau unitaire. $K$ est un corps si tout élément non nul est inversible pour $\times$. 
    Plus formellement ssi : 
        \[ \forall x \in K^*, \; \exists y \in K, \; xy = yx = 1_K \quad \text{ssi} \; \mathcal{U}(K) = K^*\] 
\end{proposition}

\begin{proposition}
    Un corps commutatif est aussi un anneau intègre. 
\end{proposition}

Tout comme pour les groupes et les anneaux, on peut définir un sous-corps d'un corps. 

\begin{definition}[Sous-corps]
    Soit $K$ un corps. Un partie $L \subseteq K$ est un sous-corps de $K$ si :
    \begin{enumerate}[label=\roman*)]
        \item $L$ est un sous-anneau de $K$
        \item $1_K \in L$ 
        \item $ \forall x \in L, x^{-1} \in L$
    \end{enumerate}
\end{definition}

En pratique pour montrer qu'un anneau est un corps, on préfèrera montrer que c'est un sous-corps quand c'est possible. 

\begin{example}
    $(\C,+,\times)$ est un corps. $(\Q,+,\times)$ et $(\R,+,\times)$ en sont des sous-corps. 
\end{example}



